% \documentclass[pra,showpacs,groupedaddress,notitlepage,superscriptaddress,floatfix,nofootinbib]{revtex4-2}
\documentclass[
    pra,
    showpacs,
    groupedaddress,
    notitlepage,
    floatfix,
    nofootinbib
]{revtex4-2}

\usepackage[T1]{fontenc}
\usepackage{lmodern,microtype}


\usepackage[english]{babel}
\usepackage{amsmath}
\usepackage{amssymb}
\usepackage{amsthm}
\usepackage{amsfonts}
\usepackage{braket}%Dirac Notation in QM
\usepackage{cases}
\usepackage{color}
\usepackage{xcolor}
\usepackage{graphicx,subfigure}% hy77 Include figure files
%\usepackage{graphics}
\usepackage{multirow}
\usepackage{placeins}
\usepackage{float}
%\usepackage{lipsum}
%\usepackage{framed}
%\usepackage{dcolumn}% Align table columns on decimal point
\usepackage{bm}% bold math
%Load hyperref last to avoid wrong equation referencing in the text
% \usepackage{hyperref}
\usepackage[hidelinks]{hyperref}

\hypersetup{
    pdftitle={Reduced Dynamics of a Qubit Coupled to a Coherent Spectator},
    pdfauthor={Jai Veilleux},
    colorlinks=true,
    linkcolor=cyan!45!black,
    citecolor=magenta!65!black,
    filecolor=magenta!65!black,
    urlcolor=cyan!45!black,
    runcolor=cyan!45!black
}
\usepackage[capitalise]{cleveref}
\usepackage{booktabs} 

\newtheorem{mytheorem}{Theorem}
\newtheorem{mylemma}{Lemma}
\newtheorem{mycorollary}{Corollary}
\newtheorem{myproposition}{Proposition}
\newtheorem{myclaim}{Claim}
\newtheorem{mydefinition}{Definition}
\newtheorem{myassumption}{Assumption}

\newcommand{\ignore}[1]{}

\newcommand{\red}[1]{\textcolor{red}{#1}}
\newcommand{\blue}[1]{\textcolor{blue}{#1}}
\newcommand{\cyan}[1]{\textcolor{cyan}{#1}}
\newcommand{\green}[1]{\textcolor{green}{#1}}
\newcommand{\magenta}[1]{\textcolor{magenta}{#1}}

\newcommand{\eps} {\varepsilon}
\newcommand{\epstot} {\eps_{\mathrm{tot}}}
\newcommand{\epserr} {\eps_{\mathrm{err}}}
\newcommand{\delerr} {\Delta_{\mathrm{err}}}
\newcommand{\therr}{\theta_{\mathrm{err}}}
\newcommand{\dphi}{\delta\varphi}
\newcommand{\dth}{\delta\theta}
\newcommand{\Herr}{H_{\mathrm{err}}}
\newcommand{\abs}[1]{\left|#1\right|}
\newcommand{\Ugrp}{\mathsf{U}}
\newcommand{\ave}[1]{\<{#1}\>}


\usepackage{commath}
\usepackage{color}
\usepackage{amsmath}
\usepackage{adjustbox}
\usepackage[normalem]{ulem}


\newcommand{\bes} {\begin{subequations}}
\newcommand{\ees} {\end{subequations}}
\newcommand{\beq}{\begin{equation}}
\newcommand{\eeq}{\end{equation}}

\def\a{\alpha}
\def\b{\beta}
\def\g{\gamma}
\def\d{\delta}
\def\e{\epsilon}
\def\ve{\varepsilon}
\def\z{\zeta}
\def\h{\eta}
\def\t{\theta}
\def\k{\kappa}
\def\m{\mu}
\def\n{\nu}
\def\x{\xi}
\def\p{\pi}
\def\r{\rho}     %%%% IMPORTANT TEx NOTE: THE REDEFINITION FOR '\r' BREAKS THE ANGSTROM SYMBOL
\def\s{XY4}
\def\ta{\tau}
\def\u{\upsilon}
\def\ph{\varphi}
\def\c{\chi}
\def\ps{\Psi}
\def\o{\omega}

\def\G{\Gamma}
\def\D{\Delta}
\def\T{\Theta}
\def\L{\Lambda}
\def\x{\xi}
\def\P{\Pi}
\def\S{XY4}
\def\U{\Upsilon}
\def\Ph{\Phi}
\def\Ps{\chi}
\def\O{\Omega}

%Some helpful quantum shortcuts
\def\ox{\otimes}
\def\>{\rangle}
\def\<{\langle}
\def\Tr{\mathrm{Tr}}
\def\Pr{\mathrm{Pr}}
\newcommand{\ketbra}[1]{|{#1}\>\!\<#1|}
\newcommand{\bracket}[1]{\<{#1}|{#1}\>}
\newcommand{\bk}[2]{\<{#1}|{#2}\>}
\newcommand{\brak}[2]{\<{#1}|{#2}\>}
\newcommand{\ketb}[2]{|{#1}\>\!\<#2|}
\newcommand{\ketbsub}[3]{|{#1}\>_{#3}\<#2|}
\newcommand{\ident}{\mathbb{}}
\newcommand{\on}[1]{\|#1\|}
\def\XY4{\text{XY}4}

\newcommand{\dd}{\,d}
\newcommand{\id}{\mathbb{I}}
\newcommand{\Ls}{\mathcal{L}}
\newcommand{\Diss}{\mathcal{D}}
\newcommand{\rhoS}{\rho_S}
\newcommand{\rhoB}{\rho_B}
\newcommand{\rhoSB}{\rho_{SB}}
\newcommand{\sx}{\sigma_x}
\newcommand{\sy}{\sigma_y}
\newcommand{\sz}{\sigma_z}
\newcommand{\spm}{\sigma_+}
\newcommand{\smm}{\sigma_-}

\newcommand{\sminus}{s_{-}}
\newcommand{\splus}{s_{+}}

\begin{document}
\title{Reduced Dynamics of a Qubit Coupled to a Coherent Spectator}
\author{Jai Veilleux}
\affiliation{University of Southern California}

\begin{abstract}
I measure the reduced dynamics of a central qubit coupled by repeated partial exchange to one spectator on IBM Quantum. 
For a spectator initialized in its ground state, the programmed circuit fixes the exact reduced-channel amplitude 
$a(n)=\cos(n\theta)$, including a phase reversal at the first return. I derive the exact Nakajima--Zwanzig memory kernel, 
the second- and fourth-order time-local predictions, and the conditions for a Markov limit. Both no-reset angle sweeps exhibit 
an increase in trace distance between the ground- and excited-input trajectories from the first ideal transfer to the first ideal 
return. In three jobs the increase is $0.753$--$0.781$ for the larger block angle and $0.618$--$0.656$ for the smaller angle. 
In the spectator-reset control, the distance decreases by $0.115$--$0.178$. Tomography of a plus-state input also resolves the 
sign change of the recovered coherence. I compare the observations with predictions set by the circuit angle and with a model 
using independently reported hardware relaxation and dephasing times. 
% No model parameters are estimated from the measured trajectories.
\end{abstract}
\maketitle

\section{Introduction}
An open system is described by tracing out environmental degrees of freedom from a joint evolution. 
An environment that retains an excitation can influence later system dynamics. A reservoir with sufficiently 
short-lived correlations can under suitable conditions yield a time-homogeneous, completely positive, trace-preserving semigroup, 
whose generator has Gorini--Kossakowski--Sudarshan--Lindblad form \cite{breuer,lidar,gorini,lindblad}. Undriven zero-temperature 
amplitude damping with a constant rate is one example. Its excited-state population falls exponentially, while the coherence falls 
at half the population rate.

A single spectator qubit gives a finite environment with an exact circuit-derived channel. As long as the spectator is retained, 
a later exchange operation can transfer an earlier excitation back to the system. Resetting the spectator instead implements a 
sequence of fresh-ancilla interactions. Collisional models of open-system dynamics and environmental memory have been demonstrated 
on IBM processors \cite{garcia}. Here I use the same exchange circuit to test both regimes on one physical qubit pair and to 
separate predictions for population, coherence, and the distinguishability of different initial states.

A return in population for one input supports the exchange picture. To test a channel property, I also compare the evolving states 
of two different preparations. A completely positive divisible family cannot recover trace distance between two states after it has 
reduced their distinguishability \cite{rivas,blp}. The experiment therefore follows ground and excited inputs across the first 
transfer and return, and uses a plus input to test the phase of the recovered coherence. These tests use fixed circuit and calibration parameters.

\section{Theory and Models}
\subsection{Partial exchange and reduced channel}
Let $S$ denote the central qubit and $B$ the spectator. Set $\hbar=1$ and $Z=\ket0\bra0-\ket1\bra1$. The free 
Hamiltonian $H_j=\omega_j\ket1_j\bra1$ places $\ket1$ at higher energy. With $\sigma_-=\ket0\bra1=(X+iY)/2$ and 
$\sigma_+=\sigma_-^\dagger$, the ideal two-qubit Hamiltonian is $H_{\mathrm{tot}}=H_S+H_B+H_{\mathrm{ex}}$, with
\begin{equation}
 H_{\mathrm{ex}}=g(\sigma_+^S\sigma_-^B+\sigma_-^S\sigma_+^B)
 =\tfrac{g}{2}(X_S X_B+Y_S Y_B).
 \label{eq:h}
\end{equation}
The ideal model starts each circuit in $\rho_{SB}(0)=\rho_S(0)\otimes(\ket0\bra0)_B$ and describes the central qubit by
$\rho_S=\Tr_B\rho_{SB}$. The spectator is the modeled environment. The calibrated model in Sec.~\ref{sec:calibrated-model}
adds local device noise to the two-qubit evolution.

Each block applies $R_{XX}(\theta)$ and then $R_{YY}(\theta)$, where 
$R_{PP}(\theta)=\exp(-i\theta P\otimes P/2)$. The generators commute. Their sum acts as $2X$ on $\{\ket{10},\ket{01}\}$ 
in system--spectator order, and vanishes on $\{\ket{00},\ket{11}\}$. For $u=n\theta$ the no-reset evolution with input $\ket{10}$ is
\begin{equation}
 U_\theta^n\ket{10}
 =\cos u\,\ket{10}-i\sin u\,\ket{01}.
 \label{eq:exchange}
\end{equation}

\begin{samepage}
For arbitrary system state $\rho$ with an uncorrelated ground-state spectator, 
the Kraus operators are $K_0=\ket0\bra0+a\ket1\bra1$ and $K_1=-i\sin u\,\ket0\bra1$, where $a=\cos u$ \cite{nielsen}. The reduced map is
\begin{equation}
 \Phi_u(\rho)=
 \begin{pmatrix}
 \rho_{00}+(1-|a|^2)\rho_{11}&a^*\rho_{01}\\
 a\rho_{10}&|a|^2\rho_{11}
 \end{pmatrix}.
 \label{eq:channel}
\end{equation}
\end{samepage}
The excited-input population is $\cos^2u$. The ground input stays fixed, so the trace distance between the 
two outputs is also $\cos^2u$. The plus input $(\ket0+\ket1)/\sqrt2$ gives $\langle X\rangle=\cos u$. At $u=\pi/2$ 
the map sends every input to the ground state. At $u=\pi$ it recovers the population and reverses the sign of the plus-state coherence. 
Since one intermediate state cannot be mapped to two distinguishable final states by a common channel, the ideal family is not 
CP divisible across that interval \cite{rivas,blp}.

\subsection{Exact memory kernels}
For the continuous Hamiltonian surrogate, use the Bohr-frequency convention
\begin{equation}
 A(t)=\sum_\omega e^{-i\omega t}A_\omega,\qquad
 A_\omega=\sum_{\epsilon_b-\epsilon_a=\omega}\Pi_a A\Pi_b,
\end{equation}
where $\Pi_a$ projects onto the energy eigenspace $\epsilon_a$ \cite{breuer,lidar}. Thus $\sigma_-^j(t)=e^{-i\omega_jt}\sigma_-^j$ and the
exchange terms oscillate at the difference frequency $\Delta=\omega_B-\omega_S$. The interaction-picture amplitudes obey
$i\dot a=g e^{-i\Delta t}b$ and $i\dot b=g e^{i\Delta t}a$. Eliminating $b$ with $b(0)=0$ gives
\begin{equation}
 \dot a(t)=-\int_0^t C(t-s)a(s)\,ds,\qquad C(t)=g^2 e^{-i\Delta t}.
 \label{eq:volterra}
\end{equation}
The free spectator ground state is stationary, so its correlation depends on the time difference. The kernel follows from
$g^2\langle\sigma_-^B(t)\sigma_+^B(0)\rangle_0=g^2e^{-i\omega_Bt}$ in the frame rotating at the system frequency.
On resonance, $\ddot a=-g^2 a$ with $a(0)=1$ and $\dot a(0)=0$, so $a(t)=\cos(gt)$. Comparing to Eq.~\eqref{eq:exchange} identifies 
the ideal angle $u=gt$. The schedule of a gate decomposition supplies a circuit duration and does not measure a continuous exchange coupling.

The coherence obeys the same convolution as $a$. The population has a separate kernel, found by taking the Laplace transform of 
$p_1(t)=\cos^2(gt)p_1(0)$:
\begin{equation}
 \widetilde f(s)=\frac{s^2+2g^2}{s(s^2+4g^2)},\qquad
 \widetilde K_p(s)=\frac{1}{\widetilde f(s)}-s
                 =\frac{2g^2s}{s^2+2g^2}.
\end{equation}
Thus $K_p(t)=2g^2\cos(\sqrt2\,gt)$. With $P_1=\ket1\bra1$ and $\Diss[L]\rho=L\rho L^\dagger-\frac12\{L^\dagger L,\rho\}$, 
the density-operator Nakajima--Zwanzig equation at resonance is
\begin{align}
 \dot\rho(t)&=\int_0^t \mathcal K(t-s)\rho(s)\,ds,\nonumber\\
 \mathcal K(t)\rho&=K_p(t)\Diss[\sigma_-]\rho
                   +[2C(t)-K_p(t)]\Diss[P_1]\rho.
 \label{eq:nz}
\end{align}
The first term propagates the population by $-K_p*p_1$, while both terms together give the coherence kernel $-C*\rho_{10}$. 
This distinguishes the operator equation from the scalar amplitude equation and keeps the memory equation regular at $a=0$ \cite{vacchini}.

\subsection{Time-local expansion from the spectator correlation}
The time-convolutionless (TCL) approach expands a time-local generator in powers of the coupling \cite{breuer,lidar}. Where $\cos(gt)\ne0$, 
the exact resonant generator is $\dot\rho=\gamma(t)\Diss[\sigma_-]\rho$ with $\gamma=-2\dot a/a=2g\tan(gt)$. The expansion at fixed $t$ gives
\begin{equation}
 \gamma_2=2g^2t,\qquad \gamma_4=2g^2t+2g^4t^3/3.
 \label{eq:rates}
\end{equation}
Integrating gives $p_1^{(2)}(u)=e^{-u^2}$ and $p_1^{(4)}(u)=e^{-u^2-u^4/6}$ for the excited input. The corresponding coherence factors are 
the positive square roots of these expressions on the near-origin branch. The derivation directly from Eq.~\eqref{eq:volterra} is in 
Appendix~\ref{app:tcl}. The exact rate has its first pole at $u=\pi/2$ and is negative on the subsequent return interval. Both truncated 
rates are positive for $t>0$ and describe the early regime. The expansion parameter is the accumulated angle $u=n\theta$. A smaller block angle
samples the early regime more finely, and both sweeps reach $u=\pi$. I score the expansions through $u=\pi/4$ using the same points for the exact channel.

\subsection{Markov limit and reset control}
The second-order cumulant master equation (C-ME) uses a coarse-graining window $T$ to average bath correlations while the system state changes
little over that window. The usual hierarchy is \cite{breuer,lidar}
\begin{equation}
 \tau_B\ll T\ll\tau_S,
 \label{eq:coarse-graining}
\end{equation}
where $\tau_B$ is the decay time of the bath correlation and $\tau_S$ is the timescale of appreciable change in the reduced state.
The window $T$ is an averaging time in the reduced theory. The block duration is set by the gate schedule. For a stationary bath
and a single decay channel, the coarse-grained population rate is
\begin{equation}
 \Gamma_T=\frac{2}{T}\mathrm{Re}\!\int_0^T (T-s)C(s)\,ds.
 \label{eq:coarse-grained-rate}
\end{equation}
When $C(s)$ decays on the scale $\tau_B$, choosing $T\gg\tau_B$ lets this rate approach $2\,\mathrm{Re}\int_0^\infty C(s)\,ds$.
For the resonant undamped spectator,
$C(s)=g^2$ for all $s$. Its correlation has no finite decay time, so Eq.~\eqref{eq:coarse-graining} has no admissible window. The Markov integral
$\int_0^\infty C(s)\,ds$ diverges and Eq.~\eqref{eq:coarse-grained-rate} gives $\Gamma_T=g^2T$.

In the secular or rotating-wave approximation to the reduced master equation, terms with distinct Bohr frequencies average out when
$|\omega-\omega'|T\gg1$. The diagonal contribution at $\omega=\omega'$ remains, so this averaging leaves the resonant correlation in
Eq.~\eqref{eq:volterra} intact \cite{breuer,lidar}. If spectator amplitude damps independently at rate $\kappa>0$, the modified correlation
$g^2e^{-(\kappa+i\Delta)t}$ has a decay time $\tau_B=1/\kappa$ and yields a population rate
\begin{equation}
 \Gamma_{\mathrm M}=2\,\mathrm{Re}\frac{g^2}{\kappa+i\Delta}
                   =\frac{2g^2\kappa}{\kappa^2+\Delta^2}
\end{equation}
in a Markov regime. Replacing $a(s)$ by $a(t)$ in Eq.~\eqref{eq:volterra} and extending the integral to infinity gives 
$\dot a=-g^2a/(\kappa+i\Delta)$. The corresponding secular amplitude-damping generator is 
$\Gamma_{\mathrm M}\Diss[\sigma_-]\rho-i[\mathrm{Im}(g^2/(\kappa+i\Delta))P_1,\rho]$. This approximation assumes slow system evolution 
over the correlation time, as when $g\ll\kappa$ on resonance. Pure spectator energy relaxation would give $\kappa=(2T_{1,B})^{-1}$ in this model. 
The hardware control uses a reset to realize memoryless exchange directly.

The ideal reset replaces the spectator by $\ket0$ after every block. The system undergoes repeated applications of $\Phi_\theta$, 
with amplitude $a_{\mathrm{reset}}(n)=(\cos\theta)^n$. This is a discrete amplitude-damping semigroup with population rate $-2\log(\cos\theta)$ 
per block for the programmed positive cosines. For $\theta=\pi/4$ the fourth-block predictions are $p_1=1/16$ for the excited input and 
$\langle X\rangle=+1/4$ for the plus input. The corresponding no-reset predictions are $p_1=1$ and $\langle X\rangle=-1$.

\section{Experimental Design}
\subsection{Circuits and job design}
I used IBM Fez with physical qubit 131 as the central qubit and 132 as the spectator. Each circuit initializes the pair, prepares the central 
qubit in $\ket0$, $\ket1$, or $\ket+$, applies a chosen number of exchange blocks and measures the central qubit in one Pauli basis. The $X$ and $Y$ 
bases use $H$ and $HS^\dagger$ immediately before computational-basis measurement \cite{tomography}. Circuits were compiled and scheduled with Qiskit 2.3.0 and
submitted through Qiskit Runtime 0.45.0 using SamplerV2 \cite{ibmsampler}. I compiled with a fixed physical layout and barriers at the
exchange boundaries. The compiled exchange-window timing agrees across preparations and measurement bases for each sweep and block count.

The $\theta=\pi/16$ sweep uses $n=0,1,2,3,4,6,8,10,12,14,16$. The $\pi/4$ no-reset and reset sweeps use $n=0,\ldots,4$. Both no-reset sweeps span the 
first transfer and return, and share nonzero angles $\pi/4,\pi/2,3\pi/4,\pi$. The reset arm resets $B$ after every exchange block, including the last. 
Each job measures all 189 distinct circuits in a different randomized order, with 512 shots per circuit. Three jobs yield 567 count records. 
The analysis treats the jobs separately.

\subsection{Calibrated block model}
\label{sec:calibrated-model}
The backend target's scheduled barriers give a duration of $0.440\,\mu\mathrm{s}$ per no-reset block and $2.124\,\mu\mathrm{s}$ per block including reset. 
The jobs share one reported calibration snapshot. For $S$, $T_1=118.2\,\mu\mathrm{s}$ and $T_2=155.2\,\mu\mathrm{s}$. For 
$B$, $T_1=99.2\,\mu\mathrm{s}$ and $T_2=112.1\,\mu\mathrm{s}$ \cite{ibmbackend}. I apply local amplitude damping and pure dephasing to both 
qubits for half each scheduled no-reset block, apply the ideal exchange of angle $\theta$, and apply noise for the other half before repeating. 
The rates are $T_1^{-1}$ and $T_2^{-1}-(2T_1)^{-1}$, so in the absence of exchange the excited population and the coherence fall at $T_1^{-1}$ and 
$T_2^{-1}$, respectively. This model uses idle calibration times within a driven gate sequence and the reset curve uses 
the ideal fresh-spectator channel.
% It does not describe detailed pulse errors, preparation, reset infidelity, or readout. 


\subsection{Tomography and uncertainty}
With ideal readout, measuring only $S$ gives the same Pauli expectations as measuring both qubits and marginalizing over $B$. The three Pauli bases determine the
central-qubit state $\hat\rho_S=(I+\hat r_X X+\hat r_Y Y+\hat r_Z Z)/2$ by linear inversion. For outcome-1 count $k_P$ in basis $P$ out of
$N$ shots, $\hat r_P=1-2k_P/N$, so $\hat p_1=(1-\hat r_Z)/2=k_Z/N$. The trace
distance for two physical qubit states is $D=\|\boldsymbol r_1-\boldsymbol r_2\|_2/2$ \cite{nielsen}. For trace-distance calculations I radially 
project a linear-inversion vector outside the unit ball onto its surface. Population comparisons use the raw $Z$ counts. Model scores are the 
mean trace distance to the predicted states and the mean absolute population error, with equal weight per preparation and sampled block count 
within the specified window.

Within each job, I independently resample each circuit's outcome-1 count from $\operatorname{Binomial}(N,k_P/N)$, repeat the reconstruction and 
projection, and use 2,000 common draws for model scores and contrasts. Intervals are pointwise 95\% percentile intervals. They describe shot 
variation conditional on each job and leave calibration uncertainty, correlated shots, and preparation and measurement errors unquantified. 
The primary contrast $\Delta D=D(u=\pi)-D(u=\pi/2)$ uses the programmed endpoints in each sweep.

\section{Results and Discussion}
\begin{figure}[!ht]
 \centering
 \includegraphics[width=\linewidth]{../results/exchange.pdf}
 \caption{(a) Central population for the excited input. 
 (b) Trace distance between the reconstructed ground- and excited-input states. 
 (c) Central $X$ expectation for the plus input. Markers are means over the three jobs and bars span the observed job range. 
 Curves are ideal circuit predictions. The vertical line marks the first ideal transfer.}
 \label{fig:hardware}
\end{figure}

\subsection{Distinguishability and phase at the first return}
The central population falls and returns in both no-reset sweeps (Fig.~\ref{fig:hardware}a). The ground- and excited-input trajectories 
approach each other at the first transfer and separate on the return. For $\theta=\pi/4$, the measured distance increases 
from $0.041$--$0.046$ at $u=\pi/2$ to $0.796$--$0.821$ at $u=\pi$ across the three jobs. The smaller-angle sweep yields the same 
qualitative pattern. In the reset control, the distance continues to fall and the excited population at the last block is $0.102$--$0.111$. 
The pointwise shot intervals for $\Delta D$ are positive for both no-reset sweeps and negative for the reset control in every job 
(Table~\ref{tab:witness}).

\begin{table}[htb]
 \centering
 \input{../results/witness.tex}
 \caption{Change $\Delta D$ from the first ideal transfer to the first return, with pointwise 95\% shot intervals, by job.}
 \label{tab:witness}
\end{table}

The plus input tests the sign of the amplitude. At $u=\pi$, $\langle X\rangle$ ranges from $-0.738$ to $-0.672$ at $\pi/16$ and from 
$-0.836$ to $-0.812$ at $\pi/4$. The reset control gives $+0.105$ to $+0.145$ (Fig.~\ref{fig:hardware}c). The negative signal in the 
no-reset sweeps is consistent with the phase reversal in Eq.~\eqref{eq:channel}. The observed distinguishability increase tests the 
predicted information return under the intended common exchange sequence. Preparation-dependent errors, readout and the different reset 
duration limit inference about a device channel in isolation. The three inputs also leave a full process reconstruction undetermined.

\subsection{Equal angles and fixed model scores}
The no-reset sweeps implement the same ideal unitary at shared angles. At $u=\pi$, the smaller-angle population is lower in all three 
jobs by $0.033$--$0.078$. Two of those three pointwise shot intervals exclude zero. At $u=\pi/2$ the smaller-angle population is higher 
in all three jobs, with one interval excluding zero. The $\pi/4$ and $3\pi/4$ contrasts have intervals spanning zero in all jobs. 
The shorter circuit uses four exchange blocks to reach $u=\pi$, compared with sixteen in the small-angle circuit. Elapsed time and gate count 
change together, leaving their effects on the contrast unresolved.

\begin{table}[htb]
 \centering
 \input{../results/scores.tex}
 \caption{Fixed predictions scored against the three measured initial states. Scores are means of the three job-level values. 
 The full window covers $0\le u\le\pi$ and the early window uses $0\le u\le\pi/4$ at the measured block counts.}
 \label{tab:scores}
\end{table}

The scheduled local-noise model yields slightly smaller mean errors across the full window for both no-reset angles (Table~\ref{tab:scores}). 
The change is modest relative to the remaining discrepancies, and the model leaves preparation and readout errors unmodeled. For example, 
the reported central-qubit readout error was about $0.118$.

Over the early window $u\le\pi/4$, TCL4 had a lower mean trace-distance error than TCL2 in each of the three jobs at both block angles.
Computed from the unrounded scores, the average reductions were about $0.0025$ at $\theta=\pi/16$ and $0.0037$ at $\theta=\pi/4$.
At both angles, the job-averaged TCL4 error was within $0.0003$ of the exact-exchange error (Table~\ref{tab:scores}).

By the first return at $u=\pi$, the exact finite-spectator channel predicts a recovery of the excited-state population and of the distinguishability
between the two inputs, along with a sign reversal of the plus-state $X$ expectation. Both no-reset sweeps exhibited all three features
(Fig.~\ref{fig:hardware}).

\section{Conclusion}
Two no-reset exchange sweeps show the return of population, the recovery of distinguishability between two initial states, and the reversal of the 
plus-state coherence. Resetting the spectator changes all three trajectories in the direction predicted by the fresh-spectator channel. 
The equal-angle measurements show a smaller population return in the longer circuit. A model fixed by the scheduled durations and measured 
idle-qubit times explains only part of the residual error. Independent readout, reset, and gate controls would be needed to resolve the 
dominant hardware contribution.

\vspace{1em}
\noindent\textit{Data and code availability.} The counts, circuit and calibration metadata, acquisition scripts, and analysis notebook 
are available at \href{https://github.com/Didgety/oqs-spec-bath}{github.com/Didgety/oqs-spec-bath}.

\begin{acknowledgments}
I thank Daniel Lidar for his lectures on open quantum systems and feedback on the project that motivated this experiment.
\end{acknowledgments}

\appendix
\section{Time-local coefficients from the memory equation}
\label{app:tcl}
Expand $a=1+a_2+a_4+\dots$ in Eq.~\eqref{eq:volterra}, where $a_j$ is order $g^j$ and vanishes initially for $j>0$. 
Substituting gives $\dot a_2=-C*1$ and $\dot a_4=-C*a_2$. At resonance,
\begin{equation}
 a_2=-\tfrac12 g^2t^2,\qquad a_4=\tfrac1{24}g^4t^4 .
\end{equation}
The logarithmic derivative through order $g^4$ is
\begin{equation}
 \frac{\dot a}{a}=\dot a_2+\dot a_4-a_2\dot a_2
                =-g^2t-\tfrac13g^4t^3,
\end{equation}
which yields Eq.~\eqref{eq:rates}. The corresponding fourth-order integrated population prediction matches the exact population through 
$u^4$. Its $u^6$ coefficient is zero, while the exact coefficient is $-2/45$. This gives a local check on the order of the approximation.

\bibliography{references}
\end{document}
